Introduction: Capital Bikeshare is a bike-share service prevalent in Washington D.C. and other metropolitan counties surrounding the area. However, as bikesharing ridership becomes common in areas with higher populations, there may be a correlation in crime rates. This can allow companies and citizens of urban areas to determine whether these services are beneficial or they can add to a growing crime rate. Our research question is does proximity to a Capital Bikeshare station correlate with higher rates of property crime but lower rates of violent crime? We decided to investigate two different types of crime to determine how Capital Bikeshare can or cannot have the potential to improve urban built environments. 

Data Summary: The data in the two datasets that we are using (Capital Bikeshare Locations and Crime Incidents in 2026) both represent populations within their parameters - the bikeshare positions of the service within the D.C metropolitan area and the population of crimes reported to the Metropolitan Police Department in Washington D.C. are represented respectively. We limited the data to strictly the DC area and locations where there are prevalent Capital Bikeshare locations. We also are focusing on incidents that are happening in 2026. This is done to standardize the dataset to minimize outliers and to see correlations in a certain area. Despite our efforts to compile accurate data, there may be some issues in our dataset. The dataset only contains reported crimes and some minor property theft crimes may be underreported, especially if they’re prevalent in certain neighborhoods. The calculations of the distances between Capital Bikeshare locations and reported crime scenes may also be inaccurate as well. These datasets work well in answering the research question because it gives a centralized location for the Capital Bikeshare locations in order to observe crime rates and the crime incident dataset also allows us to determine whether property or violent crime is more prevalent.

We have determined null and alternative hypotheses for whether the Capital Bikeshare locations have an impact on property and violent crime. The null hypothesis for property crime is There is no statistically significant correlation between a location's proximity to a Capital Bikeshare station and its rate of property crime (the correlation coefficient is ⲣ = 0). The alternate hypothesis is that there is a statistically significant correlation between proximity to a Capital Bikeshare location and a property crime incident. The null hypothesis for violent crime is There is no statistically significant correlation between a location's proximity to a Capital Bikeshare station and its rate of violent crime (the correlation coefficient is ⲣ = 0). The alternate hypothesis is that there is a statistically significant correlation between proximity to a Capital Bikeshare location and a violent crime incident.
